\chapter{从备份文件夹到分支闸门}
\label{ch:branch-pipeline}

第~\ref{ch:deadline-ai}~章讲过“能跑、正确、优化”的三段式。那时候，走这三段的是我自己。到了 Agent 时代，这三段没变，但主语变了：不再只是我走，而是我带着 Agent 一起走。

第一阶段是Make it Work。先让 Agent 把东西跑起来。塑料轮子、LL(1)、能出 C、能过 demo。这一阶段不追求优雅，只追求存在性。第二阶段是Make it right。能跑之后，不能停在“看起来对”。全量回归、parity check、边界用例、失败模式、错误处理，都要补上。AI 的价值从“生成”变成“执行”。第三阶段是Make it fast。正确之后，才谈优化。把字符串拼接换成 StringBuilder，把无序容器换成保留插入序的 OrderedMap，把模板在生成阶段内联进可执行文件。33.2 秒到 4.6 秒，不是算法变了，是写法对了。

三段不是阶梯，是循环。任何一次需求变更、硬件改版、约束变化，都可能把你打回 Work。Agent 让这个循环转得更快，但转不转、往哪转、什么时候停，仍然由你决定。原来是我一个人走 Work $\rightarrow$ Right $\rightarrow$ Fast。现在是我带着 Agent 走这三段。它负责跑，我负责停；它不知疲倦，我不知妥协。

% ---------------------------------------------------------------------
\section{git 是什么，为什么 Agent 时代你需要它}

git 是记录“谁在什么时候改了什么”的工具。提交保存一次有说明的快照，历史保存快照之间的顺序，差异告诉你具体改了哪几行，回退让你回到已知可用的状态，分支让不同任务互不干扰，合并把验收过的结果送回主线。

很多公司不用 git，原因通常很现实：项目小、人少、交付方式靠文件拷贝，老项目也没有评审流程。只要一个人慢慢改，靠备份文件夹也能撑一阵子。但 Agent 一旦进入工程，这套办法马上失效。AI 改代码速度快，一个下午可能出现几十处修改；没有 git，你不知道它改了什么、哪一版还能跑、出了问题该回到哪里。

如果公司暂时不用 git，至少在自己的项目里用起来。提交、分支、回退三个能力，就能先保住现场。不要等到 AI 把仓库改乱以后，才开始找第一份备份。

% ---------------------------------------------------------------------
\section{git 在 Agent 时代的角色变了}

以前，git 是版本控制。人写代码，人提交，历史主要用来追责和回溯。现在，git 变成了 AI 的围栏、审计日志、回滚机制和合并闸门。

围栏负责限定活动范围，AI 不直接推 \texttt{main}，只能在自己的分支上工作并提交 PR。审计日志负责留下证据，任何改动都能看 diff、看提交说明、看回归结果。回滚机制负责止损，失败就回到已知可用的提交。合并闸门负责验收，只有通过检查和人工复核的结果才能进入主线。

核心变化只有一句话：git 从“记录我改了什么”，变成“管住 AI 改了什么”。

% ---------------------------------------------------------------------
\section{我的分支模型}

我的分支模型保持简单，只分职责，不把流程画成迷宫：

\begin{itemize}
\item \texttt{main}：稳定分支，只接受验收过的合并；
\item \texttt{human/*}：人工修改分支，用于本地编译、烧录、验板和临时调整；
\item \texttt{ai/*}：AI 工作分支，用于生成、重构、补测试和跑回归；
\item \texttt{ai/regression-*}：全量回归的临时分支，跑完即删，不长期保留。
\end{itemize}

AI 在 \texttt{ai/*} 里随便折腾，\texttt{main} 永远保持干净。人工分支和 AI 分支都从 \texttt{main} 拉出，互不覆盖，最后都经过同一道 PR 质检口。

% ---------------------------------------------------------------------
\section{AI 分支的基本规则}

规则要写在仓库里，不能只靠提醒。AI 只能从 \texttt{main} 拉出 \texttt{ai/*}，不能直接修改 \texttt{main}；只能在自己的分支上提交、跑回归、提 PR；编译产物、日志和临时文件只进隔离目录，不进版本库。

\texttt{.gitignore} 要写清楚哪些目录 AI 不能碰，哪些文件不进版本库。例如构建目录、日志目录、备份目录、编辑器的临时文件，都应该提前排除。分支是权限，ignore 是边界。先画边界，再让 AI 进来干活。

% ---------------------------------------------------------------------
\section{分支流水线图}

下面这张图描述 AI 分支、人工分支、PR 质检口和 \texttt{main} 之间的关系。AI 永远不直接推 \texttt{main}，只能提 PR；人工分支和 AI 分支都从 \texttt{main} 拉出，互不干扰；合并前必须经过 PR 质检口；通过就合并，打回就继续改。

\begin{figure}[htbp]
	\centering
	\begin{tikzpicture}[
		node distance=1.1cm and 1.4cm,
		box/.style={
			rectangle,
			draw,
			rounded corners,
			align=center,
			minimum width=3.55cm,
			minimum height=1.25cm,
			font=\small
		},
		mainbox/.style={box, fill=blue!10},
		aibox/.style={box, fill=green!10},
		humanbox/.style={box, fill=orange!10},
		prbox/.style={box, fill=purple!10},
		arrow/.style={-{Latex[length=3mm]}, thick}
		]
		% 主线与两条工作分支
		\node[mainbox] (main) {\texttt{main}\\稳定分支\\只接受验收过的合并};
		\node[aibox, below left=1.5cm and 1.8cm of main] (ai)
			{\texttt{ai/*}\\AI 干活的分支\\改代码、跑全量回归、提 PR};
		\node[humanbox, below right=1.5cm and 1.8cm of main] (human)
			{\texttt{human/*}\\人工改的分支\\本地编译、烧录、验板、提 PR};

		% 统一质检口
		\node[prbox, below=2.7cm of main] (pr)
			{PR 质检口\\看 diff、回归报告、失败解释\\本地再验一遍};
		\node[mainbox, below right=1.5cm and 0.9cm of pr] (merge)
			{通过\\合并回 \texttt{main}};
		\node[aibox, below left=1.5cm and 0.9cm of pr] (reject)
			{打回\\AI 在 \texttt{ai/*} 继续改};

		% 分支流转
		\draw[arrow] (main) -- (ai);
		\draw[arrow] (main) -- (human);
		\draw[arrow] (ai) -- (pr);
		\draw[arrow] (human) -- (pr);
		\draw[arrow] (pr) -- (merge);
		\draw[arrow] (pr) -- (reject);
		\draw[arrow] (reject) -- (ai);
	\end{tikzpicture}
	\caption{Agent 时代的分支流水线：AI 是车间，PR 是质检口，\texttt{main} 是出厂库}
	\label{fig:branch-pipeline}
\end{figure}

把这张图落到工程里，只需要记住一句：AI 是车间，PR 是质检口，\texttt{main} 是出厂库。

% ---------------------------------------------------------------------
\section{PR 作为质检口}

PR 不是形式，也不是给领导看的截图。它是代码进入稳定分支之前的最后一道检查。先看 diff，确认改动范围与任务一致；再看回归报告，确认全量测试、边界用例和失败解释；最后卡三类问题：有没有动不该动的目录，有没有夹带无关改动，关键用例有没有人工复核。

审查完之后，我会本地拉下 \texttt{ai/*}，在自己的编译环境里再验一遍。没问题，合并；有问题，打回，AI 在原分支上继续改。AI 可以提交很多次，但只有经过质检的结果才算完成。

% ---------------------------------------------------------------------
\section{编译目录隔离与止损}

编译产物、日志、临时文件只进隔离目录，不进 git，也不进源码树。隔离目录让 AI 可以反复构建、反复清理，不会把生成物和源代码混在一起。还要设置止损。连续若干轮回归失败后，AI 必须停下并报告，不许继续猜着改。止损规则写进 \texttt{AGENTS.md}，不能靠口头约定，也不能靠人的耐心临时判断。这是在保护你的现场。

% ---------------------------------------------------------------------
\section{关键用例的最终解释权}

AI 可以跑全量回归，输出通过和失败。但“通过”能不能信，边界用例要不要人工复核，最终由人决定。特别是涉及时钟、内存布局、中断优先级、协议兼容和高风险硬件操作的用例，必须由人确认结果的含义。

这不是不信任自动化，而是承认自动化只能回答“有没有触发条件”，不能替你回答“这个结果在项目里算不算正确”。AI 可以跑全量，但通过的解释权在人。

% ---------------------------------------------------------------------
\section{全量备份与 git 的分工}

全量备份是大改之前的保险丝，放在 git 之外，AI 不可写。它的作用不是日常版本管理，而是防止仓库本身被误删、误覆盖或彻底破坏。git 分支和 PR 则是日常工作的主刹车，负责记录改了什么、怎么回退、怎么合并。

两者不冲突：备份防“AI 把仓库搞烂”，git 管“改了什么、怎么回、怎么合”。备份是安全气囊，git 是刹车系统。两个都在，才敢让 Agent 提速。

% ---------------------------------------------------------------------
\section{从个人技巧到可复制工作流}

一开始，我的经验很像个人技巧：中午自费、低峰错峰、窗口伪装，尽量不影响别人，也尽量不被别人打断。这些办法能救急，但带不走。换一个环境、换一台电脑、换一个同事，就全部失效。

后来我把它们固化成分支模型、PR 闸门、止损规则和 \texttt{AGENTS.md}。前者是个人套利，后者是团队可以照着执行的流程。个人套利是战术，制度化工作流是资产。战术帮你活下来，资产帮你把下一次也做好。

% ---------------------------------------------------------------------
\section{给团队的最小落地清单}

先不要追求完整平台，把最小闭环跑起来：

\begin{itemize}
\item 建好 \texttt{main}、\texttt{human/*}、\texttt{ai/*} 三条分支线；
\item 写 \texttt{.gitignore}，排除编译输出、日志、临时文件和备份目录；
\item 写 \texttt{AGENTS.md}，写明分支规则、PR 规则、止损规则和禁止触碰的目录；
\item 规定 AI 只能提 PR，不能直接推 \texttt{main}；
\item 关键用例人工复核，通过的解释权留在人手里。
\end{itemize}

先用起来，再优化；先有围栏，再放 AI 进去。等团队习惯了这套流程，再讨论更复杂的自动化，顺序才不会反。
